\chapter{从采集样本生成字形}\label{ch:generation}
\entry{统一应用顶部\ui{生成}标签页。默认选择\ui{v6.2 exp（推荐）}；另可选择\ui{v6.1 exp}、\ui{v4 adaptive}。}

\section{准备作者与样本}
填写与采集文件相同的作者 ID，选择该作者的汉字单字 JSON 目录。在生成页直接选择样本目录；采集输出目录可在采集页状态栏查看。目录中应有真实采集的 \code{handwriting_single_v1} 文件；生成轨迹、字体包、编辑字形不能代替真实采集样本建档。数字、字母及数学符号通过字体库直接加入排版资源，不进入本页的汉字部件建档流程。

\section{导入、建档、生成}
\begin{enumerate}
\item 点\ui{导入样本}。程序核对作者、轨迹和动态字段，显示接收数量及跳过数量；详细原因保存于导入记录。没有有效样本时，不能继续建档。
\item 选择生成方式，点\ui{建立档案}。v6.1 自动准备完整来源轨迹及标准笔画映射；v4 准备自适应风格档案。无需另建副档案。部分样本因笔序或映射不可靠而不能成为生成来源，建档记录会列出原因。
\item 在\ui{生成内容}输入汉字，或勾选\ui{生成结构库全部汉字}。点\ui{开始生成}，查看进度。可以取消；串行在当前字结束后停止，多进程停止投递并回收本任务工作进程，保留已核验提交的字。
\item 在生成页点击\ui{制作字体…}，在弹出的制作窗口中为结果命名，作者与生成结果已经带入；按需要加入同一作者的符号样本。制作成功后回到字体库并选中新字体，再选择\ui{用于排版}。
\end{enumerate}
未满足条件的操作按钮会禁用。缺少导入索引时会提示先导入；切换方法后，没有对应档案就需要建档。正在操作时作者和方法暂时锁定。采集目录新增文件后，重新导入再建档；应用不会悄悄更新已生成的结果。

\note{每次导入、建档和生成均保存为独立批次。再次点击生成会建立新批次，不覆盖旧样本、档案或结果。本版全库任务支持取消并保留完成部分，尚未提供应用内断点续跑入口。}

\section{生成方法}
\textbf{v6.2 exp} 在 v6.1 基础上增加竖弯钩与止部关键笔语义准入，拒绝不兼容来源后重新检索，保持真实笔迹及原几何检查。未找到来源时仍明确提示参考替代；来源完整不等于人工质量合格。user\_001 验证合并字体仅重生孔、正，孔的末笔保留参考缺口，其余字体候选和旧文档保持原样。
\textbf{v6.1 exp} 使用经过验证的参考修正、完整部件候选筛选、禁止拆分与位置变体规则、受限关键笔替换及作者单笔回退。保留作者来源点列、供体、变换、拒绝原因和缺口。关键笔替换仅限当前已验证的羊部件第六笔规则，不是任意笔画的自动修复。程序从本次档案重新寻找供体，不会载入另一作者的历史生成字冒充新结果。

本版沿用跨字生成规则：目标字本身的样本不作供体，暂不启用同字原稿优先复用。历史人工冻结的供体和布局只适用于相应档案，不会自动套用到新作者。来源覆盖率不能直接解释为允许同字复用时的生产覆盖率，也不等于人工评价合格率。

\textbf{v4 adaptive} 使用固定的自适应部件融合、强度 1.0 和原有生成路径。它没有完整可靠的逐笔作者来源证明，字体中按来源未知提示。应用不提供神经生成、训练或 CUDA 检测入口。两种方法均不承诺重复运行能产生经过验证的自然随机变体。

详细规则可悬停生成方式旁的\ui{i}提示查看，无需调整部件强度等内部参数。

\section{来源缺口、失败与复核}
作者来源不足但参考几何有效的字可以生成和排版，明确标记参考替代笔序及原因。字体应用时按文档实际使用的字汇总警告；允许继续使用，不强制立即补采。任务完成、作者来源完整、协议有效和人工评价是不同状态。非有限坐标、损坏轨迹或来源校验失败的字不会被当作成功字发布。

生成页显示简短错误和诊断记录位置，不直接显示异常堆栈。部分字失败时，结果目录保留失败记录，字体只纳入完成的字；取消后未生成部分不被当成完成。发布字体前应核对完成数与请求数。

保存为 \code{.inkdoc} 会冻结实际候选、轨迹、度量和来源状态。重开不会悄悄重新生成或换字。补采与重生成后，应制作新字体并重新检查；旧评价不自动证明新结果合格。

\section{记录与复现}
工作区的 \code{generation} 目录分别保存导入、档案、结果与诊断。v6.1 的结果含 \code{protocol.json}、逐字轨迹压缩文件和来源侧记录；v4 保留轨迹目录，另有实际方法及固定配置记录。开发者字体保留生成依赖、配置、版本与哈希；用户字体保留必要版本、来源摘要和警告。完整记录不表示已打包所有原始数据，独立复现仍需要相应样本与结构资源。

排版后如需查看来源缺口，在排版页打开\ui{设置 → 问题 → 字体提示}。主窗口不再单独显示跨模块工具栏；字体库直接通过同名标签页打开。

\section{字体库与生成页的分工}
生成页按准备样本、生成汉字、生成结果组织操作。运行记录默认折叠，需要查看详细路径和处理原因时展开。字体库默认只展示已经制作或导入的字体，可以用于排版、导出用户版或验证。

从字体库点击\ui{制作字体…}，也可选择历史生成结果或仅添加符号样本。制作窗口不会调用生成器；它将已有成品轨迹封装为字体。制作过程中输入锁定，成功后自动选中新字体。跨作者交接时不会沿用另一作者的符号目录。

\section{按字集与多进程（0.1.40）}
\ui{内置简繁字表…}提供大陆《通用规范汉字表》2013 年版 8105 字、GB/T 12345-1990 繁体汉字 6866 字、台湾《常用國字標準字體表》1982 年版 4808 字，可自由合并。另有按 Unihan 17.0.0 全对应构造的 8593 字派生繁体扩展集，不属于地区规范。显示来源与版本，合并去重保留共用字；一简多繁、自身对应及地区编码异体保留独立身份，不进行用户文本转换。

结构库缺字完整列出，可复制报告。只有明确勾选已查看缺口并仅生成支持部分后，才允许排除缺口；改变选择须重新确认。完整原表、实际目标、缺口、版本与哈希均随任务保存。选择字表不改变供体池，也不保证同一编码的地区笔形；本版没有香港规范完整字表。台湾源表 A02149 的编码栏冲突有明确勘误及核对依据。

\ui{导入字表…}仍支持 UTF-8 TXT 或 CSV；多列 CSV 必须有唯一的 \code{character}、汉字或字符列。去重保留首次顺序，不做简繁转换，不拆开变体序列。自定义字表不支持的字符仍阻止任务启动。编辑内置字表后改为自定义字表，原字表身份不再沿用。

\ui{从文档提取汉字…}支持 Markdown 和已有项目；\ui{从当前排版提取}读取当前未保存源码及模板配置，不改原文、设置或冻结选样。读取正文、公式、代码和图题中的汉字，不读 URL 或图片路径；当前模板只纳入活动布局引用的字段。没有模板上下文的 Markdown 采用保守文本集合，原解析诊断保留。提取结果仍可编辑。

\ui{计算方式}提供自动、串行、二进程和四进程。自动模式少于 128 字使用串行，大任务同时受可用内存和 CPU 限制，上限四进程；v4 始终串行。按字分派，每个进程复用完整作者档案，不缩小供体池，不省略语义、碰撞或来源检查。

v6.1/v6.2 任务创建独立 SQLite 一致性备份，包含已提交 WAL 内容；逻辑身份只在此不可变资源视图内复用。新任务重新读取，不按路径或修改时间跨任务缓存。\code{target-manifest.json} 保存实际字符与字表身份；\code{protocol.json} 保存规则、资源、计算方式、失败、取消及待完成项。

工作进程不可用时，记录原因并串行补齐未提交字符；已有完成字不重算。只承认逐字完整提交与哈希一致的结果，残留暂存文件不等于完成。取消后只保留已提交结果；尚未提供自动断点续跑。制作字体仍另存新包，旧候选、旧文档与冻结选样不自动迁移。
